% Selected books on stochastic processes in ecology and epidemiology
% Requires natbib if the parent document uses \citep commands.
\section{Selected Books on Stochastic Processes in Ecology and Epidemiology}
\label{sec:selected_stochastic_books}

Stochastic approaches to ecological population dynamics and infectious diseases have a long tradition. The following four books provide a deliberately short selection rather than an exhaustive bibliography. Together, they trace the development of the field from early stochastic population and epidemic models to modern algorithms for simulating individual events in large populations and networks. They also complement the seven research papers discussed in the preceding section.

\begin{enumerate}
\item \textbf{M. S. Bartlett (1960).} \emph{Stochastic Population Models in Ecology and Epidemiology}. Methuen, London (also distributed by Wiley, New York). \citep{Bartlett1960StochasticPopulation}

This short monograph is a pioneering introduction to stochastic models of ecological populations and infectious diseases. It discusses birth--death processes, extinction, population interactions, epidemic models and recurrent outbreaks, emphasizing aspects of population dynamics that cannot be understood from deterministic equations alone. In particular, Bartlett's work helped establish demographic stochasticity, epidemic fade-out and the persistence of infection as central questions in mathematical epidemiology.

\item \textbf{R. M. Nisbet and W. S. C. Gurney (1982).} \emph{Modelling Fluctuating Populations}. John Wiley \& Sons, Chichester. \citep{NisbetGurney1982Fluctuating}

This book places population fluctuations at the centre of ecological modelling. Rather than focusing exclusively on equilibria and stability, it explains how deterministic dynamics, demographic noise and varying environments contribute to observed temporal variability. It develops mathematical tools for analysing fluctuations, oscillations and extinction, with applications to interacting populations and spatially heterogeneous environments. It provides a particularly useful bridge between theoretical population ecology and the stochastic methods introduced in this course.

\item \textbf{N. G. van Kampen (2007, 3rd edition).} \emph{Stochastic Processes in Physics and Chemistry}. Elsevier, Amsterdam. \citep{vanKampen2007Stochastic}

A standard mathematical reference on stochastic processes, applicable well beyond physics and chemistry. It treats Markov processes, master equations, one-step processes, stationary distributions, first-passage questions, Fokker--Planck equations and systematic approximations such as the system-size expansion. It is especially valuable for students wishing to move beyond the elementary derivations presented in the lectures and develop a more rigorous understanding of the relationship between microscopic stochastic events and macroscopic dynamics. The first edition appeared in 1981.

\item \textbf{N. Masuda and C. L. Vestergaard (2022).} \emph{Gillespie Algorithms for Stochastic Multiagent Dynamics in Populations and Networks}. Cambridge University Press (Cambridge Elements). \citep{MasudaVestergaard2022Gillespie}

This concise and openly accessible volume provides a modern introduction to exact simulation of continuous-time stochastic systems with interacting agents. It explains the Direct and First Reaction methods, efficient implementations including the Next Reaction Method, and extensions to systems on networks and to non-Poissonian dynamics. Its emphasis on practical simulation makes it a natural companion to the second day of the course, particularly for students interested in spatially extended ecological or epidemic models. The full text is available at \url{https://doi.org/10.1017/9781009239158}.
\end{enumerate}

\noindent\textbf{A historical note.} The references are listed in the proposed reading order, using the third edition (2007) for van Kampen. Strictly by \emph{first publication}, van Kampen's book (1981) precedes Nisbet and Gurney (1982). Together, these books illustrate both the long history of stochastic thinking in ecology and epidemiology and its continuing development.

% Companion to selected_books_section.tex.
% Requires natbib (or biblatex with natbib compatibility) and the Allen2010 entry.
\section{An Introductory Textbook for Further Study}
\label{sec:introductory_stochastic_textbook}

The four books listed above provide a historical and methodological perspective on the development of stochastic modelling in ecology and epidemiology. For students approaching the subject without extensive mathematical training, however, a more introductory textbook may be a useful starting point.

\noindent\textbf{L. J. S. Allen (2010).} \emph{An Introduction to Stochastic Processes with Applications to Biology}, 2nd edition. Chapman \& Hall/CRC. \citep{Allen2010}

This textbook introduces stochastic processes through biological examples rather than developing the mathematics in isolation. Starting from basic probability concepts, it covers discrete- and continuous-time Markov chains, birth--death processes, population extinction and models of infectious disease dynamics. Its worked examples and exercises help students connect transition rates, probability distributions and stochastic simulations with biological questions. Although some familiarity with calculus and elementary probability is helpful, it is a particularly suitable entry point for biology and environmental-science students wishing to study the ideas of this course in greater depth before consulting more mathematically demanding references such as van Kampen.

